% Options for packages loaded elsewhere
\PassOptionsToPackage{unicode}{hyperref}
\PassOptionsToPackage{hyphens}{url}
\PassOptionsToPackage{dvipsnames,svgnames,x11names}{xcolor}
%
\documentclass[
  11pt,
]{article}
\usepackage{amsmath,amssymb}
\usepackage{iftex}
\ifPDFTeX
  \usepackage[T1]{fontenc}
  \usepackage[utf8]{inputenc}
  \usepackage{textcomp} % provide euro and other symbols
\else % if luatex or xetex
  \usepackage{unicode-math} % this also loads fontspec
  \defaultfontfeatures{Scale=MatchLowercase}
  \defaultfontfeatures[\rmfamily]{Ligatures=TeX,Scale=1}
\fi
\usepackage{lmodern}
\ifPDFTeX\else
  % xetex/luatex font selection
\fi
% Use upquote if available, for straight quotes in verbatim environments
\IfFileExists{upquote.sty}{\usepackage{upquote}}{}
\IfFileExists{microtype.sty}{% use microtype if available
  \usepackage[]{microtype}
  \UseMicrotypeSet[protrusion]{basicmath} % disable protrusion for tt fonts
}{}
\makeatletter
\@ifundefined{KOMAClassName}{% if non-KOMA class
  \IfFileExists{parskip.sty}{%
    \usepackage{parskip}
  }{% else
    \setlength{\parindent}{0pt}
    \setlength{\parskip}{6pt plus 2pt minus 1pt}}
}{% if KOMA class
  \KOMAoptions{parskip=half}}
\makeatother
\usepackage{xcolor}
\usepackage[margin=1in]{geometry}
\usepackage{longtable,booktabs,array}
\usepackage{calc} % for calculating minipage widths
% Correct order of tables after \paragraph or \subparagraph
\usepackage{etoolbox}
\makeatletter
\patchcmd\longtable{\par}{\if@noskipsec\mbox{}\fi\par}{}{}
\makeatother
% Allow footnotes in longtable head/foot
\IfFileExists{footnotehyper.sty}{\usepackage{footnotehyper}}{\usepackage{footnote}}
\makesavenoteenv{longtable}
\setlength{\emergencystretch}{3em} % prevent overfull lines
\providecommand{\tightlist}{%
  \setlength{\itemsep}{0pt}\setlength{\parskip}{0pt}}
\setcounter{secnumdepth}{-\maxdimen} % remove section numbering
\ifLuaTeX
  \usepackage{selnolig}  % disable illegal ligatures
\fi
\IfFileExists{bookmark.sty}{\usepackage{bookmark}}{\usepackage{hyperref}}
\IfFileExists{xurl.sty}{\usepackage{xurl}}{} % add URL line breaks if available
\urlstyle{same}
\hypersetup{
  pdftitle={Research paths from the GR / Mercury-perihelion papers --- a prioritized response},
  pdfauthor={Trey Morris with Claude Opus 4.7},
  pdfsubject={Five candidate research paths opened by Zorn's 2026-06-04
papers, ordered by priority for the dual-theory framework},
  colorlinks=true,
  linkcolor={blue},
  filecolor={Maroon},
  citecolor={Blue},
  urlcolor={blue},
  pdfcreator={LaTeX via pandoc}}

\title{Research paths from the GR / Mercury-perihelion papers --- a
prioritized response}
\author{Trey Morris with Claude Opus 4.7}
\date{2026-06-04}

\begin{document}
\maketitle

\hypertarget{research-paths-from-the-gr-mercury-perihelion-papers-for-zorn}{%
\section{Research paths from the GR / Mercury-perihelion papers --- for
Zorn}\label{research-paths-from-the-gr-mercury-perihelion-papers-for-zorn}}

\textbf{To:} Zorn (@Zorns-Lemmon). \textbf{From:} Trey Morris (with
Claude Opus 4.7). \textbf{Date:} 2026-06-04. \textbf{Re:} Your
2026-06-04 email\footnote{Email from Zorn (@Zorns-Lemmon), 2026-06-04,
  preserved verbatim in the first comment on issue \#139.} on issue
\#139\footnote{PyPhysics issue \#139, ``Parse Zorn's GR/Mercury papers
  and ship a research-paths response PDF to Zorn.''}. Thank you for the
weekend hunt. The papers you sent open five distinct threads. This note
orders them by how directly they bear on the dual-theory framework, and
gives a concrete next step for each.

\begin{center}\rule{0.5\linewidth}{0.5pt}\end{center}

\hypertarget{cover-note}{%
\subsection{Cover note}\label{cover-note}}

This is a planning response, not a derivation. It ranks the threads your
papers open and proposes a sequence. It does not yet contain results
from those papers. The Bootello\footnote{A. Bootello, \emph{Journal of
  Modern Physics} (SCIRP), paper id 145485. Published PDF and HTML at
  the journal site.} and Wex\footnote{N. Wex, \emph{Testing Relativistic
  Gravity with Radio Pulsars}, arXiv:1402.5594.} sources are not yet
parsed into the repository, and the Edvardsson\footnote{S. Edvardsson,
  \emph{Relativistic gravitational force}, \emph{Celestial Mechanics and
  Dynamical Astronomy} (no arXiv copy located; Lab subscription
  unavailable as of 2026-06-04).} paper is access-blocked at the Lab.
Where a paragraph interprets a source's relevance to the framework, that
interpretation is flagged for your review.

Per the lab's writing standard\footnote{\texttt{Roadmapping/Tooling/VOICE\_MATCH\_GILL.md},
  §3.bis--§3.quinquies (equation, sentence, citation, and emoji
  discipline).} and AI-use compliance\footnote{\texttt{Roadmapping/Tooling/CROCCO\_COMPLIANCE.md}.},
this report is substantive AI end-to-end. Every ranking and every
framework-mapping claim is an interpretive move and is marked for your
sign-off. No specific result, equation, or number is attributed to a
paper we have not yet read. The standard textbook results quoted below
as benchmarks are footnoted to the repository's existing perihelion
documents, not to fabricated external citations.

\begin{center}\rule{0.5\linewidth}{0.5pt}\end{center}

\hypertarget{the-ordering-at-a-glance}{%
\subsection{§1. The ordering, at a
glance}\label{the-ordering-at-a-glance}}

You asked, in effect, which of these is worth pursuing. The order we
propose is \textbf{1, 2, 4, 3, 5}, read as priority for the framework,
not as the order you raised them.

\begin{longtable}[]{@{}
  >{\raggedright\arraybackslash}p{(\columnwidth - 8\tabcolsep) * \real{0.2000}}
  >{\raggedright\arraybackslash}p{(\columnwidth - 8\tabcolsep) * \real{0.2000}}
  >{\raggedright\arraybackslash}p{(\columnwidth - 8\tabcolsep) * \real{0.2000}}
  >{\raggedright\arraybackslash}p{(\columnwidth - 8\tabcolsep) * \real{0.2000}}
  >{\raggedright\arraybackslash}p{(\columnwidth - 8\tabcolsep) * \real{0.2000}}@{}}
\toprule\noalign{}
\begin{minipage}[b]{\linewidth}\raggedright
Rank
\end{minipage} & \begin{minipage}[b]{\linewidth}\raggedright
Path
\end{minipage} & \begin{minipage}[b]{\linewidth}\raggedright
Source
\end{minipage} & \begin{minipage}[b]{\linewidth}\raggedright
Why here
\end{minipage} & \begin{minipage}[b]{\linewidth}\raggedright
Status
\end{minipage} \\
\midrule\noalign{}
\endhead
\bottomrule\noalign{}
\endlastfoot
1 & Mercury perihelion via Bootello & Bootello\footnote{A. Bootello,
  \emph{Journal of Modern Physics} (SCIRP), paper id 145485. Published
  PDF and HTML at the journal site.} & Plugs straight into the live
perihelion thread\footnote{\texttt{Roadmapping/Author\_Reports/2026-05\_corda\_perihelion\_review\_for\_gill.md}
  --- end-to-end review of the Corda framework; source of the \(+37.8\)
  vs \(-7.2\) arcsec/century adjudication.}\footnote{\texttt{Roadmapping/Author\_Reports/2026-05\_perihelion\_classical\_derivation\_for\_gill.md}
  --- standard-GR and proper-time perihelion derivations via Binet's
  equation; source of the benchmark (1), the residual figure, and the
  identity (2).} & \textbf{Lead} \\
2 & Interaction-based gravity & Edvardsson\footnote{S. Edvardsson,
  \emph{Relativistic gravitational force}, \emph{Celestial Mechanics and
  Dynamical Astronomy} (no arXiv copy located; Lab subscription
  unavailable as of 2026-06-04).} & Deepest conceptual fit to the
proper-time framing & \textbf{High-upside, access-blocked} \\
3 & Pulsar-timing proper-time tests & Wex\footnote{N. Wex, \emph{Testing
  Relativistic Gravity with Radio Pulsars}, arXiv:1402.5594.} & A real
experimental channel for a proper-time prediction &
\textbf{Experimental} \\
4 & Gravitoelectromagnetism & (your GEM reading) & Connects your weekend
reading to our EM work & \textbf{Structural} \\
5 & Light-bending consistency & shared & A cross-check gate, not a
standalone program & \textbf{Validation} \\
\end{longtable}

The reasoning for each rank follows.

\begin{center}\rule{0.5\linewidth}{0.5pt}\end{center}

\hypertarget{path-1-mercury-perihelion-via-bootello-lead}{%
\subsection{§2. Path 1 --- Mercury perihelion via Bootello
(Lead)}\label{path-1-mercury-perihelion-via-bootello-lead}}

This is the lead because it lands on top of work already in motion. The
repository already carries a Mercury-perihelion derivation built on
Corda's framework\footnote{C. Corda, \emph{The secret of planets'
  perihelion between Newton and Einstein}, \emph{Physics of the Dark
  Universe} \textbf{32} (2021) 100834.}, reviewed end to end\footnote{\texttt{Roadmapping/Author\_Reports/2026-05\_corda\_perihelion\_review\_for\_gill.md}
  --- end-to-end review of the Corda framework; source of the \(+37.8\)
  vs \(-7.2\) arcsec/century adjudication.}, and worked through Binet's
orbit equation both the standard-GR way and the proper-time
way\footnote{\texttt{Roadmapping/Author\_Reports/2026-05\_perihelion\_classical\_derivation\_for\_gill.md}
  --- standard-GR and proper-time perihelion derivations via Binet's
  equation; source of the benchmark (1), the residual figure, and the
  identity (2).}. The benchmark is the general-relativistic anomalous
advance per orbit,

\[
\Delta\phi_{\text{GR}} \;=\; \frac{6\pi G M}{c^{2}\,a\,(1-e^{2})}, \tag{1}
\]

which accounts for the observed Mercury residual of roughly forty-three
arcseconds per century\footnote{\texttt{Roadmapping/Author\_Reports/2026-05\_perihelion\_classical\_derivation\_for\_gill.md}
  --- standard-GR and proper-time perihelion derivations via Binet's
  equation; source of the benchmark (1), the residual figure, and the
  identity (2).}.

The reason a second independent derivation matters is that our own
thread is not yet closed. The framework's proper-time update sits in an
unresolved adjudication. One route gives a Corda-style result near
\(+37.8\) arcseconds per century; the proper apsidal-angle route gives
near \(-7.2\) arcseconds per century\footnote{\texttt{Roadmapping/Author\_Reports/2026-05\_corda\_perihelion\_review\_for\_gill.md}
  --- end-to-end review of the Corda framework; source of the \(+37.8\)
  vs \(-7.2\) arcsec/century adjudication.}. The structural identity
behind that tension is

\[
\frac{\Delta\phi_{\text{framework}}}{\Delta\phi_{\text{GR}}} \;=\; -\frac{1}{6}, \tag{2}
\]

which falls out of the algebra as an identity, not an
approximation\footnote{\texttt{Roadmapping/Author\_Reports/2026-05\_perihelion\_classical\_derivation\_for\_gill.md}
  --- standard-GR and proper-time perihelion derivations via Binet's
  equation; source of the benchmark (1), the residual figure, and the
  identity (2).}.

So Bootello is useful precisely as an outside check. If we parse it into
the same Binet apparatus, we can ask whether his derivation corroborates
one of our two framework numbers, reproduces the GR value (1), or stands
apart from all three. A third independent perihelion calculation is the
cheapest way to break the current tie.

\textbf{Next step.} Download the published Bootello PDF\footnote{A.
  Bootello, \emph{Journal of Modern Physics} (SCIRP), paper id 145485.
  Published PDF and HTML at the journal site.} into
\texttt{External\_Papers/}, parse it to markdown, and slot its force law
into Binet's equation alongside the Corda and proper-time chains. One
table, three derivations, same apparatus.

\begin{center}\rule{0.5\linewidth}{0.5pt}\end{center}

\hypertarget{path-2-edvardssons-interaction-based-gravity-high-upside-access-blocked}{%
\subsection{§3. Path 2 --- Edvardsson's interaction-based gravity
(High-upside,
access-blocked)}\label{path-2-edvardssons-interaction-based-gravity-high-upside-access-blocked}}

This is ranked second because it is the deepest conceptual fit, gated
only by access. You relayed that the abstract gives up the principle of
equivalence and rebuilds gravity as a traditional interaction force
inside special relativity, while claiming consistency with the Mercury
precession and light bending\footnote{S. Edvardsson, \emph{Relativistic
  gravitational force}, \emph{Celestial Mechanics and Dynamical
  Astronomy} (no arXiv copy located; Lab subscription unavailable as of
  2026-06-04).}. That is structurally close to the dual theory's stance.
The dual theory's distinguishing feature is its proper-time,
interaction-first formulation of dynamics. A gravity model that is
interaction-first and equivalence-principle-free is the same shape of
idea applied to gravity.

The upside is high because, if the abstract holds, Edvardsson is a
worked example of the framework's own philosophy in the one regime ---
gravitation --- where the dual theory has so far said the least. The
risk is that we cannot yet read it. It is in \emph{Celestial Mechanics
and Dynamical Astronomy}, the Lab has no subscription, and there is no
arXiv copy\footnote{S. Edvardsson, \emph{Relativistic gravitational
  force}, \emph{Celestial Mechanics and Dynamical Astronomy} (no arXiv
  copy located; Lab subscription unavailable as of 2026-06-04).}.

\textbf{Next step.} Try three acquisition routes in order: an author
preprint or personal-page copy, an institutional or interlibrary-loan
request, then a direct note to the author. Until we have the full text,
we do not derive anything from it. We hold the conceptual parallel as a
hypothesis, not a result.

\begin{center}\rule{0.5\linewidth}{0.5pt}\end{center}

\hypertarget{path-4-pulsar-timing-proper-time-tests-experimental}{%
\subsection{§4. Path 4 --- Pulsar-timing proper-time tests
(Experimental)}\label{path-4-pulsar-timing-proper-time-tests-experimental}}

You flagged this as fertile ground, and we agree it is the strongest
\emph{experimental} channel. Binary-pulsar timing measures the
periastron advance directly, among other post-Keplerian effects, at high
precision\footnote{N. Wex, \emph{Testing Relativistic Gravity with Radio
  Pulsars}, arXiv:1402.5594.}. The perihelion question of Path 1 and the
precision-timing question here are the same physics observed in two
systems: Mercury in the solar field, and a neutron star in a companion's
field.

The framework's contribution would be a predicted fractional correction
to the periastron advance,

\[
\dot{\omega}_{\text{obs}} \;=\; \dot{\omega}_{\text{GR}}\,\bigl(1 + \delta_{\text{framework}}\bigr), \tag{3}
\]

where \(\delta_{\text{framework}}\) is the proper-time correction the
dual theory predicts. The same \(-1/6\)-style structural factor that
appears in (2) for Mercury should, if the framework is consistent,
reappear in \(\delta_{\text{framework}}\) here. A binary pulsar is where
that prediction would actually be tested against data, since pulsar
timing reaches a precision the solar-system perihelion does not.

\textbf{Next step.} Pull the Wex review's LaTeX source\footnote{N. Wex,
  \emph{Testing Relativistic Gravity with Radio Pulsars},
  arXiv:1402.5594.}, identify the binary systems with the cleanest
measured periastron advance, and fold a \(\delta_{\text{framework}}\)
forecast into the experimental-design plan\footnote{\texttt{Roadmapping/Experimental\_Design\_Plan\_Dual\_Theory.md}.}.
This is a longer thread than Paths 1 and 2, which is why it sits behind
them.

\begin{center}\rule{0.5\linewidth}{0.5pt}\end{center}

\hypertarget{path-3-gravitoelectromagnetism-structural}{%
\subsection{§5. Path 3 --- Gravitoelectromagnetism
(Structural)}\label{path-3-gravitoelectromagnetism-structural}}

This connects your weekend reading to the repository's electromagnetism
work\footnote{\texttt{Roadmapping/Electromagnetism/} --- the framework's
  electromagnetism verification and derivation work.}. In the
weak-field, slow-motion limit, linearized general relativity can be
written in a Maxwell-analog form, with a gravito-electric field and a
gravito-magnetic field obeying field equations of the schematic shape

\[
\nabla\cdot\mathbf{E}_g \;=\; -4\pi G\,\rho, \qquad
\nabla\times\mathbf{B}_g \;=\; -\frac{4\pi G}{c^{2}}\,\mathbf{j} + \frac{1}{c^{2}}\frac{\partial \mathbf{E}_g}{\partial t}. \tag{4}
\]

The convention-dependent numerical factors in (4) are left for the human
pass to pin against a canonical reference, since they vary by a factor
of up to four across the literature.

The reason this is structural rather than a lead is that it is a
reformulation route, not a new prediction by itself. The interesting
question is whether the dual theory's existing EM structure maps onto
the gravito-magnetic sector cleanly, and whether that map reproduces the
perihelion advance of Path 1 by a different road. If it does, GEM
becomes a unifying language for Paths 1 and 4. If it does not, it tells
us where the framework's EM-gravity analogy breaks.

\textbf{Next step.} Rewrite the GEM field equations (4) in the dual
theory's own variables and check whether the gravito-magnetic term
yields the same periastron advance the Binet route gives in Path 1.

\begin{center}\rule{0.5\linewidth}{0.5pt}\end{center}

\hypertarget{path-5-light-bending-consistency-validation}{%
\subsection{§6. Path 5 --- Light-bending consistency
(Validation)}\label{path-5-light-bending-consistency-validation}}

This is ranked last on purpose. It is a gate, not a program. Both
Edvardsson\footnote{S. Edvardsson, \emph{Relativistic gravitational
  force}, \emph{Celestial Mechanics and Dynamical Astronomy} (no arXiv
  copy located; Lab subscription unavailable as of 2026-06-04).} and any
GEM route\footnote{\texttt{Roadmapping/Electromagnetism/} --- the
  framework's electromagnetism verification and derivation work.} claim
to reproduce gravitational light deflection, so the standard benchmark

\[
\alpha \;=\; \frac{4 G M}{c^{2}\,b}, \tag{5}
\]

with \(b\) the impact parameter, is the shared cross-check every
candidate mechanism must pass. It is most useful applied to the outputs
of Paths 2 and 3, rather than pursued on its own.

\textbf{Next step.} Once any candidate mechanism produces a deflection
prediction, compare it against (5). Treat agreement as a necessary
condition, not as a finding in itself.

\begin{center}\rule{0.5\linewidth}{0.5pt}\end{center}

\hypertarget{proposed-sequence-and-what-we-need-from-you}{%
\subsection{§7. Proposed sequence and what we need from
you}\label{proposed-sequence-and-what-we-need-from-you}}

The short version. Start Path 1 now, because it sharpens a tie we are
already trying to break. Start the \emph{acquisition} for Path 2 in
parallel, because the bottleneck there is access, not effort. Treat
Paths 4 and 3 as the next wave, and use Path 5 as the gate they must
clear.

Two requests. First, the second arXiv link you mislaid\footnote{Email
  from Zorn (@Zorns-Lemmon), 2026-06-04, preserved verbatim in the first
  comment on issue \#139.} --- when you find it, drop it on issue \#139
and we will fold it into the same parse. Second, if you have any back
channel to the Edvardsson paper\footnote{S. Edvardsson,
  \emph{Relativistic gravitational force}, \emph{Celestial Mechanics and
  Dynamical Astronomy} (no arXiv copy located; Lab subscription
  unavailable as of 2026-06-04).}, that unblocks the highest-upside
thread.

\begin{center}\rule{0.5\linewidth}{0.5pt}\end{center}

\hypertarget{references}{%
\subsection{References}\label{references}}

\end{document}
